\documentclass[11pt, a4paper]{article}
\usepackage{xeCJK}

\usepackage[top=0.7in, bottom=0.7in, left=0.65in, right=0.65in]{geometry}
\usepackage{titlesec}
\usepackage{enumitem}
\usepackage{hyperref}
\usepackage{fontspec}
\usepackage{tabularx}
\setmainfont{texgyretermes}[
    Extension = .otf,
    UprightFont = *-regular,
    BoldFont = *-bold,
    ItalicFont = *-italic,
    BoldItalicFont = *-bolditalic
]


% 设置段落间距
\setlength{\parindent}{0pt}
\setlength{\parskip}{6pt}

% 自定义命令
\newcommand{\namesection}[2]{
    \begin{center}
        {\Huge \textbf{#1}} \\
        \vspace{5pt}
        {\large #2}
    \end{center}
    \vspace{10pt}
}

\titleformat{\section}{\large\bfseries}{}{0em}{}[\titlerule]
\titlespacing{\section}{0pt}{10pt}{5pt}

\titleformat{\subsection}{\bfseries}{}{0em}{}
\titlespacing{\subsection}{0pt}{8pt}{2pt}

\newenvironment{tightemize}{
    \begin{itemize}[leftmargin=*, nosep]
    \setlength{\itemsep}{3pt}
    \setlength{\parskip}{0pt}
}{\end{itemize}}

\begin{document}

\namesection{李师贤个人简历}{ 
\href{mailto:zxlishixian@gmail.com}{zxlishixian@gmail.com}  |   +86-16673588836
}

\section{个人概述}

\textbf{\hspace{2em}北京邮电大学未来学院（元班）2023级本科生，主修计算机科学与技术方向，并接受电子信息、通信、人工智能等多学科交叉培养。数学基础扎实，数学分析 90、高等代数 93、概率论与随机过程 97，获全国大学生数学竞赛国家级二等奖。曾赴台湾阳明交通大学资讯工程学系交换学习，修读编译器、数据库、深度学习、计算机视觉等课程并取得全 A 及以上成绩，具备英文课程学习、论文阅读与跨文化交流经验。研究兴趣集中于人工智能、多智能体系统、计算机视觉与智能系统，期待在科研训练中进一步发挥数学建模、跨学科整合与工程实现能力。}

\section{教育背景}

\subsection{北京邮电大学}
\textbf{未来学院（元班） | 2023级本科生} \\

\textbf{绩点：3.58/4  |  专业课成绩：87/100} \\
专业排名：20/60

\textbf{核心课程与专业基础:}
\begin{tightemize}
    \item \textbf{数学基础}: 数学分析 (90)、高等代数(93)、概率论与随机过程 (97)
    \item \textbf{计算机与人工智能基础}: 数据结构与算法(91)、离散数学(92)、编译器设计(A)、数据库系统(A)、深度学习(A)
    \item \textbf{信息通信与电子基础}: 信号与系统(91)、通信原理(86)、数字信号处理(90)
    \item \textbf{网络与安全方向}: 计算机网络(87)、信息网络导论(95)
\end{tightemize}

\vspace{0.2cm}

\subsection{台湾阳明交通大学（交换经历）}
\textbf{資訊工程學系 | 2025年秋季交换生} \\
\textbf{研修课程成绩单} \\

\begin{tabularx}{\textwidth}{@{}>{\centering\arraybackslash}p{0.15\textwidth}X>{\centering\arraybackslash}p{0.12\textwidth}>{\centering\arraybackslash}p{0.14\textwidth}@{}}
\textbf{課程性質} & \textbf{課程名稱} & \textbf{學分} & \textbf{等級成績} \\
\\
  選 & 編譯器設計概論          & 3.00 & A \\
  選 & 資料庫系統概論\#        & 3.00 & A \\
  研 & 電腦視覺實務與深度學習  & 3.00 & A \\
  研 & 深度學習                & 3.00 & A \\
\end{tabularx}

\vspace{0.1cm}
{\footnotesize  \textbf{说明：}研：大學部修研究所課程；\#：英語授課}

{\small \textbf{课程工程实践:} 交换期间完成覆盖编译器、数据库系统、深度学习与计算机视觉方向的多项课程设计，包括完整编译器流程实现、校园 Lost and Found 数据库系统及前端开发、深度学习课程竞赛、计算机视觉大作业与顶会论文复现改进等。相关实践积累了较好的工程实现、系统设计、实验分析与技术报告撰写能力，也进一步训练了面向科研问题的复现、比较和改进意识。部分代码与材料见\href{https://github.com/zxlishixian?tab=repositories}{GitHub repositories}。}

\vspace{0.2cm}
\section{项目经历}

\subsection{北京中医药大学“岐黄智教”大创项目}
\textbf{核心成员} | 获北京市"挑战杯"人工智能专项赛道国家三等奖
\begin{tightemize}
    \item 项目旨在运用人工智能技术解决传统中医教学中的难点
    \item 作为核心成员，主要负责项目\textbf{网站搭建与模型接入}工作
    \item 项目成果荣获人工智能专项赛道国家级三等奖，并持续产出成果
\end{tightemize}

\vspace{0.22cm}

\subsection{北京邮电大学“智创助手”大创项目}
\textbf{核心成员}
\begin{tightemize}
    \item 开发可\textbf{本地化部署}的分析工具使用手册自然语言模型
    \item 负责\textbf{本地一键化部署与封装}等技术研究
\end{tightemize}

\vspace{0.22cm}

\subsection{北京邮电大学“暑期日本电气通信大学短期交流PBL”项目}
\textbf{小车自动巡航项目组组长} | 获得优秀结项证书
\begin{tightemize}
\item 基于 ROS 和 Gazebo 搭建仿真环境，在虚拟机上完成小车模型、激光雷达与场景的模拟测试
\item 编写 Python 节点实现键盘遥控、激光雷达 SLAM 建图及实时地图可视化
\item 集成 ROS Navigation Stack，实现一键设定目标点后的全局路径规划与动态避障
\item 将仿真代码部署至实体小车，调试传感器驱动并验证自动巡航与避障功能
\end{tightemize}

\vspace{0.22cm}

\subsection*{低光照图像增强中的自适应强度坍缩研究}
\textbf{TAICA台灣大專院校人工智慧學聯盟專區《電腦視覺實務與深度學習》课程结课设计}
\begin{tightemize}
    \item 复现\textbf{CVPR论文“HVI: A New Color Space for Low-light Image Enhancement”}，发现其强度坍缩机制因全局固定参数\(k\)难以兼顾不同曝光图像的噪声抑制与细节保留。
    \item 提出\textbf{内容感知自适应坍缩策略}：基于图像全局平均亮度动态调整坍缩强度，暗区强坍缩以抑制噪声、亮区弱坍缩以保留细节，且不修改原始CIDNet架构。
    \item 在LOL-v1数据集上验证：改进方法在训练早期（前200轮）PSNR提升1.33 dB、SSIM提升0.04、LPIPS降低0.033，暗区噪声减少、边缘结构更稳定。
    \item 撰写技术报告，通过定量指标与可视化结果证实表示层优化可有效提升低光照增强的收敛速度与初期质量。
\end{tightemize}

\vspace{0.22cm}

\subsection{校级“雏雁计划”TrafficFlow项目}
\textbf{核心成员} | 获北京邮电大学校级一等奖
\begin{tightemize}
    \item 基于深度学习的信号灯时长优化系统，提高城市交通流畅性
    \item 主要负责\textbf{前期技术调研与数据预处理}工作
\end{tightemize}

\section{获奖与证书}

\begin{tightemize}
    \item \textbf{全国大学生数学竞赛} - 国家级二等奖
    \item \textbf{北京市"挑战杯"大学生创新创业竞赛} - “人工智能+”专项赛道 国家级三等奖
    \item \textbf{雏雁计划TrafficFlow项目} - 校级一等奖
    \item \textbf{鸿雁杯创新创业大赛} - 校级三等奖
    \item \textbf{北京邮电大学校级二等奖学金} (2023-2024学年)
    \item \textbf{北京邮电大学校级二等奖学金} (2024-2025学年)
    \item \textbf{校级优秀团员} (2023-2024学年)
    \item \textbf{英语能力}: 雅思 6.5分 | 大学英语六级 (CET-6)  538分 | 大学英语四级 (CET-4)  609分 ，具备熟练的英文文献阅读能力
    \item \textbf{暑期PBL项目结项证书}: 圆满完成小车自动巡航项目，获得优秀结项证书
    \item \textbf{交换生修业证书}: 圆满完成2025秋季学期台湾地区大学交换生学业并取得所有课程全A及以上成绩，获得荣誉证书
\end{tightemize}

\end{document}
